% !TEX root = ../../main.tex
% C5.1 — Thiết kế kiến trúc hệ thống (RÚT GỌN 2026-09-29: bỏ 5.1.2 "Hai luồng xử lý chính" vì trùng 5.2; bản cũ ở report/archive/)
% Nguồn: architect.md mục 3 (ranh giới thành phần, ba nơi lưu trữ), mục 1 (tuần tự, concurrency 1), mục 9, 13 (#4 RAM: chưa cần tách encoder).
% Khớp code: InProcessQueryEncoder trong app.py (encoder truy vấn chạy trong tiến trình API, trừ khi đặt PERSON_SEARCH_ENCODER_URL);
% hàng đợi = bảng processing_job trong PostgreSQL, worker nhận việc bằng lease + heartbeat (workers/durable.py);
% Vite proxy /api → Flask (frontend/vite.config). Hình H2: figures/H2-kien-truc.drawio (tools/h2_architecture_drawio.py).
\section{Thiết kế kiến trúc hệ thống}
\label{sec:5.1}

Hình~\ref{fig:architecture} trình bày kiến trúc tổng thể của hệ thống, tổ chức quanh hai giai đoạn đã nêu ở Mục~\ref{sec:3.1.6}: luồng \emph{lập chỉ mục} ở cột trái biến dữ liệu camera thành các lần xuất hiện có thể tìm kiếm, và luồng \emph{tìm kiếm} ở cột phải phục vụ thao tác của người dùng. Hai luồng gặp nhau ở tầng lưu trữ phía dưới.

\begin{figure}[!tbp]
\centering
% Rộng 18 cm, vượt đều mỗi bên lề 1 cm cho dễ đọc (lề trái còn 2 cm, lề phải còn 1 cm).
\makebox[\textwidth][c]{\includegraphics[width=18cm]{H2-kien-truc.png}}
\caption{Kiến trúc tổng thể của hệ thống}
\label{fig:architecture}
\end{figure}

\subsection{Các thành phần của hệ thống}
\label{sec:5.1.1}

\begin{itemize}
    \item \textbf{Nguồn dữ liệu camera.} Các tệp video của bộ dữ liệu thử nghiệm được FFmpeg phát liên tục lên máy chủ MediaMTX thành các luồng RTSP; mỗi luồng là một camera logic (Mục~\ref{sec:3.7.3}) và đối với hệ thống không khác camera IP thật.
    \item \textbf{Ứng dụng web.} Giao diện xây dựng bằng React, dùng chung cho cả ba vai trò; ứng dụng web chỉ giao tiếp với máy chủ ứng dụng, không truy cập trực tiếp kho dữ liệu nào.
    \item \textbf{Máy chủ ứng dụng.} Một tiến trình Flask cung cấp API cho toàn bộ chức năng, và nạp sẵn bộ mã hóa RaSa để mã hóa truy vấn tìm kiếm.
    \item \textbf{Tiến trình xử lý nền.} Một tiến trình Python riêng thực hiện toàn bộ phần xử lý AI: tự tạo các phiên xử lý RTSP xoay vòng giữa các camera, lần lượt nhận từng công việc trong hàng đợi và chạy quy trình phân tích.
    \item \textbf{Tầng lưu trữ.} PostgreSQL lưu dữ liệu nghiệp vụ và hàng đợi công việc; Milvus lưu vector đặc trưng kèm các trường dùng để lọc; MinIO lưu khung hình toàn cảnh đại diện của mỗi lần xuất hiện.
\end{itemize}

\subsection{Các nguyên tắc thiết kế}
\label{sec:5.1.3}

Nguyên tắc quan trọng nhất là tách xử lý AI khỏi máy chủ ứng dụng. Trên máy triển khai, việc phân tích video chậm hơn thời gian thực khoảng bảy lần, nên máy chủ ứng dụng không đọc luồng RTSP hay chạy quy trình phân tích mà chỉ tạo công việc. Tiến trình xử lý nền thực hiện các công việc này một cách độc lập. Người dùng vẫn tìm kiếm bình thường trong khi hệ thống đang lập chỉ mục, và lỗi khi phân tích không làm dừng máy chủ ứng dụng.

Hàng đợi công việc được đặt ngay trong PostgreSQL: mỗi công việc là một bản ghi, thay vì dùng thêm một hệ thống hàng đợi thông điệp riêng. Cách này giảm số thành phần phải triển khai trên máy cá nhân và giữ trạng thái công việc cùng chỗ với dữ liệu nghiệp vụ. Tiến trình xử lý nền thực hiện tuần tự một công việc tại một thời điểm, còn các phiên RTSP có thời lượng giới hạn được xoay vòng giữa các camera, để mọi camera đều lần lượt được phân tích mà không làm quá tải máy chỉ có CPU.

Về lưu trữ, mỗi kho giữ một loại dữ liệu và PostgreSQL là nơi quyết định. Ba phần của cùng một lần xuất hiện được liên kết bằng một mã chung, và mọi kết quả do Milvus trả về đều được đối chiếu lại với PostgreSQL trước khi hiển thị. Mọi truy cập dữ liệu, kể cả ảnh, đều đi qua máy chủ ứng dụng, để ma trận phân quyền ở Mục~\ref{sec:4.6} được áp dụng tại một điểm duy nhất.

Cuối cùng, hai luồng dùng chung bộ mã hóa và quy trình. Truy vấn được mã hóa bằng đúng trọng số RaSa đã dùng khi lập chỉ mục, để hai bên nằm trong cùng một không gian vector. Luồng RTSP và tệp video chỉ khác nhau ở bước đọc khung hình.
